% Generated draft: requires mathematical and source review.
\documentclass[11pt]{article}
\usepackage[margin=25mm]{geometry}
\usepackage{fontspec}
\setmainfont{Latin Modern Roman}
\usepackage{amsmath,amssymb,mathtools,unicode-math}
\setmathfont{Latin Modern Math}
\usepackage{enumitem,needspace,xurl,hyperref}
\setlength{\parindent}{0pt}
\setlength{\parskip}{5pt}
\setlength{\emergencystretch}{3em}
\newcommand{\N}{\mathbb{N}}
\newcommand{\Z}{\mathbb{Z}}
\newcommand{\Q}{\mathbb{Q}}
\newcommand{\R}{\mathbb{R}}
\newcommand{\C}{\mathbb{C}}
\newcommand{\F}{\mathbb{F}}
\newcommand{\A}{\mathbb{A}}
\newcommand{\PP}{\mathbb{P}}
\newcommand{\E}{\mathbb{E}}
\newcommand{\eps}{\varepsilon}
\newcommand{\dd}{\,\mathrm{d}}
\begin{document}
\section*{UnsolvedMath ID: 30002224}
\textbf{Source problem code:} catalogue entry
\subsection{Definitions and mathematical statement}
For a unital complex \(C^*\)-algebra \(A\), let \(K_0(A)\) be the Grothendieck group of stable Murray--von Neumann equivalence classes of projections in matrix algebras over \(A\), and let \(K_1(A)\) be the group of homotopy classes of stabilized unitary matrices over \(A\).

The term \emph{easy quantum group} refers to a compact matrix quantum group whose intertwiner spaces are spanned by linear maps associated with specified set partitions. More explicitly, in the usual uncoloured orthogonal convention, let \(u\) be the fundamental representation on \(\C^n\). For a partition \(p\) of \(k\) upper and \(\ell\) lower points, define
\[
T_p(e_{i_1}\otimes\cdots\otimes e_{i_k})
=\sum_{j_1,\ldots,j_\ell=1}^{n}
\delta_p(i,j)\,e_{j_1}\otimes\cdots\otimes e_{j_\ell},
\]
where \(\delta_p(i,j)=1\) if indices belonging to the same block of \(p\) are equal, and is zero otherwise. Easiness means that, for a suitable category of partitions \(\mathcal P\),
\[
\operatorname{Hom}_G(u^{\otimes k},u^{\otimes\ell})
=\operatorname{span}\{T_p:p\in\mathcal P(k,\ell)\}
\qquad(k,\ell\geq0),
\]
with zeroth tensor power equal to the trivial representation. Coloured unitary conventions also involve conjugate representations. The supplied record does not specify which convention or class of partition categories is intended; the displayed definition explains one standard convention without resolving that scope ambiguity.

\textbf{Problem, with unresolved scope.} For a general easy quantum group \(G\), determine the two operator \(K\)-theory groups
\[
K_0(A(G))\qquad\text{and}\qquad K_1(A(G)),
\]
where \(A(G)\) denotes the associated \(C^*\)-algebra intended by the source. The request is a general computation problem, rather than a conjectured formula supplied by the record. The record does not specify whether \(A(G)\) is the universal function algebra \(C^u(G)\), the reduced function algebra \(C^r(G)\), or another associated algebra. These choices must be resolved before the question has a unique precise target; no equality of their \(K\)-theory groups is assumed. Nor does the record specify whether ordered or unital \(K_0\)-data are part of the requested computation.
\subsection{Short English statement}
Determine \(K_0\) and \(K_1\) of the \(C^*\)-algebras associated with easy quantum groups in general. The intended class of easy quantum groups and the associated algebra or completion require clarification. The supplied curation labels the general problem open and reports computations for important subclasses, including free quantum groups and fixed-point algebras; these are supplied literature assessments, not independently verified conclusions.
\subsection{Sources}
\begin{itemize}
\item Local UnsolvedMath registry: supplied complete record with permanent ID \(30002224\), exact source problem code \texttt{catalogue entry}, in the Oberwolfach Reports Open Problems collection. No individual proposer is identified. The registry gives proposed year \(2012\), status \texttt{open}, and statement status \texttt{corrected\_verified}; these are supplied metadata, not independent mathematical verification.
\item Supplied report citation: \emph{C*-Algebras, Dynamics, and Classification} (2013), Oberwolfach Reports, pp.\ 3129--3209, DOI \texttt{10.4171/owr/2012/52}. \url{https://doi.org/10.4171/owr/2012/52}. The supplied citation year is retained despite the 2012 proposed year and DOI component.
\item Supplied literature source, labelled ``Oberwolfach report, original open problem'': \url{https://publications.mfo.de/bitstream/handle/mfo/3323/OWR_2012_52.pdf?isAllowed=y&sequence=1}.
\item Supplied literature source, labelled ``Vergnioux--Voigt, K-theory of free quantum groups'': \url{https://doi.org/10.1007/s00208-013-0902-9}.
\item Supplied literature source, labelled ``Gabriel--Weber, fixed-point algebras'': \url{https://doi.org/10.3842/SIGMA.2016.097}.
\item Literature assessment dated exactly \texttt{2026-08-21}: the supplied curation labels the problem open, reports computable \(K\)-theory for important subclasses including free quantum groups and fixed-point algebras, and describes the general question as unresolved. This is supplied dated curation derived from the OWR record, not independently verified source consultation or literature assessment. Time-sensitive status and cited claims require independent verification.
\end{itemize}
\end{document}
